% ID: SEL-D0005-OA-E04
% ESTADO: validada. S006, 2026-11-07; tres pasadas del mismo autor y revisión matemática, didáctica y visual.
\textbf{Modelo y estimador.}\par\nopagebreak[4]
La proporción muestral es $\widehat p=16/60=4/15$, que se conserva sin redondear. Con muestreo aleatorio independiente usamos el intervalo normal aproximado para la proporción; los recuentos 16 y 44 permiten la aproximación usual. No se conoce tamaño poblacional para una corrección finita.

\textbf{a) Confianza del 90\,\%.}\par\nopagebreak[4]
$\alpha=0.10$, luego $z_{0.95}\approx1.644853627$ (tabla interpolada: 1.645). El intervalo es
\[
\widehat p\pm z_{0.95}\sqrt{\frac{\widehat p(1-\widehat p)}{60}}.
\]
\[
\boxed{IC_{\text{90\%}}\approx(0.172762100,\ 0.360571233).}
\]
Equivale aproximadamente a entre 17.28\,\% y 36.06\,\%. La confianza se refiere a la cobertura del procedimiento para la proporción poblacional en muestras repetidas.

\textbf{b) Tamaño mínimo para error inferior a 0.1.}\par\nopagebreak[4]
Con el mismo cuantil y la misma proporción piloto,
\[
z_{0.95}\sqrt{\frac{(4/15)(11/15)}n}<0.1
\iff n>\frac{z_{0.95}^2(4/15)(11/15)}{0.1^2}
\approx52.908405.
\]
El tamaño total mínimo es \boxed{n=53\text{ empleados}}. Con 52 el error supera 0.1; con 53 es menor. También se obtiene 53 usando 1.645. Que el tamaño requerido sea inferior a 60 es coherente: la muestra inicial ya cumplía esta precisión. No se pide añadir 53 empleados ni se debe usar 0.1 como amplitud total del intervalo.
